% answers_short.tex —— 2016 年·全国Ⅰ卷·文科 · 简短答案速查
% 由 tools/make_short_answers.py 自 answers.tex 的 \gksolhead/\gkanswer 提取，
% 并与 2016/2016_数学真题_全国1文_答案要点.md 逐题比对后生成，勿手改（改源 MD 后重跑）。

\begin{gknotice}
\textbf{使用说明}：本节只给最终答案便于快速核对，完整推导见解析版。答案为编写者独立推导，\textbf{非官方标准答案}。
\end{gknotice}

\gkpart{选择题}
\begin{center}\begin{tabular}{ccccccccccccc}
\toprule
题号 & 1 & 2 & 3 & 4 & 5 & 6 & 7 & 8 & 9 & 10 & 11 & 12 \\
\midrule
答案 & B & A & C & D & B & D & A & B & D & C & A & C \\
\bottomrule
\end{tabular}\end{center}

\gkpart{填空题}
\begin{itemize}[leftmargin=2.6em]
  \item[\textbf{13.}] \(-\frac23\)
  \item[\textbf{14.}] \(-\frac43\)
  \item[\textbf{15.}] \(4\pi\)
  \item[\textbf{16.}] \(216000\) 元
\end{itemize}

\gkpart{解答题}
\begin{itemize}[leftmargin=2.6em]
  \item[\textbf{17.}] (1) \(a_n=3n-1\)；(2) \(S_n=\frac32(1-3^{-n})\)
  \item[\textbf{18.}] (1) \(G\) 为 \(AB\) 中点；(2) \(F\) 是 \(E\) 在平面 \(PAC\) 上的正投影，\(V_{PDEF}=\frac43\)
  \item[\textbf{19.}] (1) $y=\begin{cases}3800,&x\le19,\\500x-5700,&x>19\end{cases}$；(2) \(n_{\min}=19\)；(3) 购买 19 个时平均费用为 \(4000\) 元，购买 20 个时为 \(4050\) 元；应购买 19 个
  \item[\textbf{20.}] (1) \(\frac{|OH|}{|ON|}=2\)；(2) 除 \(H\) 外没有其他公共点
  \item[\textbf{21.}] (1) 当 \(a\ge0\) 时在 \(x=1\) 左递减、右递增；当 \(-\frac e2<a<0\) 时令 \(\xi=\ln(-2a)<1\)，递增、递减、递增区间依次为 \((-\infty,\xi),(\xi,1),(1,+\infty)\)；当 \(a=-\frac e2\) 时在 \(\mathbb R\) 上严格递增；当 \(a<-\frac e2\) 时令 \(\xi=\ln(-2a)>1\)，递增、递减、递增区间依次为 \((-\infty,1),(1,\xi),(\xi,+\infty)\)；(2) \(a>0\)
\end{itemize}

\gkpart{选做题}
\begin{itemize}[leftmargin=2.6em]
  \item[\textbf{22.}] (1) \(AB\) 与圆相切；(2) \(AB\parallel CD\)
  \item[\textbf{23.}] (1) 圆心 \((0,1)\)、半径 \(a\) 的圆；\(\rho^2-2\rho\sin\theta+1-a^2=0\)；(2) U；若另加 \(C_1\cap C_2\ne\varnothing\)，则 \(a=1\)。见 ISSUE-0932
  \item[\textbf{24.}] (1) $f(x)=\begin{cases}x-4,&x<-1,\\3x-2,&-1\le x<\frac32,\\4-x,&x\ge\frac32\end{cases}$；(2) \((-\infty,\frac13)\cup(1,3)\cup(5,+\infty)\)
\end{itemize}
